\section{Emptiness and the Two Truths}

This is the conceptual base of Mahayana and Tibetan Buddhism. It is the view of reality that underlies every list of realms, paths, and stages in the tradition. It is not itself a list. It is the model the lists hang on.

\subsection*{Emptiness}

\textbf{Emptiness} (\textbf{shunyata}) is the central concept. It means that phenomena lack \textbf{inherent existence}. Nothing exists from its own side, on its own, by its own essence. The thing emptiness denies is \textbf{own-being} (\textbf{svabhava}), the imagined property of standing alone, of being self-contained, of carrying a fixed independent identity. Emptiness is the absence of own-being.

A thing is empty because it depends. It depends on causes. It depends on parts. It depends on the mind that labels it. A chariot is not found in any one of its parts, nor anywhere apart from them. The word chariot is laid onto a collection. So the chariot is empty of being a chariot from its own side.

What emptiness amounts to can be put as three balances.

\begin{itemize}
 \item Empty of inherent existence, not empty of existence.
 \item Empty of a fixed essence, not empty of function.
 \item Empty of standing alone, not empty of appearing.
\end{itemize}

\subsection*{What emptiness does not mean}

Emptiness is not nothingness. This is the most important guardrail in the system. To say a flower is empty is not to say there is no flower. It is to say the flower has no independent, self-standing core. The flower still grows, still has color, still can be picked. It functions. It simply does not exist in the solid way it appears to.

Nagarjuna warned that emptiness misunderstood as nothingness is like grasping a snake by the wrong end. It harms the one who holds it. The view steers between two errors.

\begin{itemize}
 \item \textbf{Eternalism}, the error of reifying, of thinking things have a solid essence.
 \item \textbf{Nihilism}, the error of denying, of thinking nothing exists or nothing matters.
\end{itemize}

Emptiness is the middle. It lets things be real conventionally while denying them ultimate solidity.

\subsection*{Dependent origination, the reason for emptiness}

Things are empty because they arise in dependence. \textbf{Dependent origination} (\textbf{pratityasamutpada}) is the reason for emptiness. The logic is tight. If a thing arises depending on other things, it cannot also be independent. Dependence and inherent existence cannot both be true. So dependence proves emptiness. Whatever is dependently arisen is what is meant by empty.

Three levels of dependence are taught.

\begin{itemize}
 \item \textbf{Causal dependence}. A sprout depends on a seed.
 \item \textbf{Dependence on parts}. A whole depends on its parts.
 \item \textbf{Conceptual dependence}. A thing depends on being labeled by mind.
\end{itemize}

The deepest level is the third. Even time, space, and the self exist only as designations laid on a basis. This is the subtle dependent origination stressed by the Prasangika Madhyamaka.

\subsection*{The two truths}

Reality is read at two levels at once. This is the doctrine of the \textbf{two truths}.

\textbf{Conventional truth} (\textbf{samvriti satya}) is how things appear and function in ordinary experience. Cause and effect, day and night, you and me. All valid at this level. The word also carries the sense of a concealing veil, since the conventional surface hides the deeper nature.

\textbf{Ultimate truth} (\textbf{paramartha satya}) is the way things actually are. Empty of inherent existence. Found by analysis, not by ordinary perception.

The two are not two separate worlds. They are two truths about one reality. The conventional flower and the empty flower are the same flower seen two ways. This is why a buddha can act in the world while seeing its emptiness.

\subsection*{The middle way and the four-cornered logic}

The school built on these ideas is \textbf{Madhyamaka}, the \textbf{Middle Way}, founded by Nagarjuna around the second century. It is the highest of the four tenet systems in the Tibetan curriculum. The middle is between two extremes. Not the extreme of existence, which is eternalism. Not the extreme of non-existence, which is nihilism. Things are like illusions, dreams, and reflections. Apparent yet empty.

Madhyamaka uses a four-cornered logic called the \textbf{tetralemma} (\textbf{catuskoti}). It tests a claim against four positions and rejects all four as ways of grasping inherent existence. Applied to whether things arise, the four corners are these.

\begin{itemize}
 \item Not from self.
 \item Not from other.
 \item Not from both.
 \item Not from neither, that is, not without any cause.
\end{itemize}

By negating all four, the analysis refuses to land on any reified position. The point is not a hidden fifth answer. The point is to exhaust grasping and let the mind rest in non-conceptual openness.

\subsection*{The four tenet systems}

The Tibetan curriculum ranks Buddhist philosophy into four schools, two lower and two higher. Each is studied as a rung. The lower views are stepping-stones to the higher ones.

\begin{longtable}{@{} c L{2.6cm} L{2.0cm} L{6.4cm}@{}}
\caption{The four Buddhist tenet systems, climbing from realism to full emptiness.}\\
\toprule
\# & School & Vehicle & Core claim about reality \\
\midrule
1 & Vaibhashika & Lower & External atoms and momentary mind-instants are truly, substantially real, and directly perceived. \\
2 & Sautrantika & Lower & External objects exist but are known only through mental representations, not directly. \\
3 & Cittamatra & Higher & There are no external objects. Only mind is real. The all-ground projects the apparent world. \\
4 & Madhyamaka & Higher & Nothing has intrinsic existence. All phenomena are empty, even mind itself. \\
\bottomrule
\end{longtable}

The fourth school splits into two branches over how emptiness is established. \textbf{Svatantrika}, associated with Bhaviveka, uses autonomous syllogisms and accepts that things exist conventionally by their own characteristics. \textbf{Prasangika}, associated with Buddhapalita and Chandrakirti, uses only consequence arguments that draw out the absurdity in the opponent's view, and it denies inherent existence even conventionally. The Gelug school, following Tsongkhapa, holds Prasangika as the highest and most refined view.

\subsection*{The three marks of existence}

Underneath all of this sits the basic Buddhist reading of reality, shared with the earlier tradition. These are the \textbf{three marks of existence}.

\begin{itemize}
 \item \textbf{Impermanence} (\textbf{anitya}). All conditioned things arise and pass. Nothing is static. There is change from moment to moment.
 \item \textbf{Suffering} or unsatisfactoriness (\textbf{duhkha}). Conditioned existence cannot give lasting satisfaction. Grasping at the impermanent brings pain.
 \item \textbf{No-self} (\textbf{anatman}). There is no permanent, independent, unitary self in persons or in things.
\end{itemize}

No-self is the bridge to emptiness. The early teaching denied a fixed self in the person. The Mahayana extended the denial to all phenomena. So Tibetan Buddhism speaks of two selflessnesses, the selflessness of persons and the selflessness of phenomena. Emptiness is the full, generalized form of no-self. It is the same insight, scaled to everything that exists. From the three marks the whole structure of view, path, and fruition unfolds.
